\documentclass[10pt]{article}
\usepackage[margin=20mm]{geometry}
\usepackage{amsmath,amssymb,mathtools,fontspec,unicode-math,xurl,hyperref}
\IfFontExistsTF{DejaVu Serif}{\setmainfont{DejaVu Serif}}{\setmainfont{Latin Modern Roman}}
\setmathfont{latinmodern-math.otf}
\setlength{\parindent}{0pt}\setlength{\parskip}{6pt}\sloppy
\title{M25-0120
Scientific Review of
Rotation-invariant web bases from hourglass plabic graphs}
\author{}
\date{26 October 2026}
\begin{document}
\maketitle
\section*{1. Identity, scope, and assessment}
Authors. Christian Gaetz, Oliver Pechenik, Stephan Pfannerer, Jessica Striker, and Joshua P. Swanson.\par
Reviewed source. The preserved 84-page author manuscript arXiv:2306.12501v6 (https://arxiv.org/abs/2306.12501v6), submitted 8 May 2026; its title page is dated 12 May 2026. The exact PDF SHA256 is\par
00000TOKEN\par
The journal record is Inventiones Mathematicae 243, 703–804 (2026), DOI 10.1007/s00222-025-01385-2 (https://doi.org/10.1007/s00222-025-01385-2), first published online 10 November 2025. The journal PDF was not obtained. This review makes no assertion of equivalence to that version. All source page numbers below refer to the preserved manuscript, whose printed and PDF page numbers agree.\par
Assessment. Two substantive proof defects are confirmed by explicit finite graph constructions. The first concerns the characterization of full reducedness by boundary-started trip strands. The second concerns the immediate-decrease assertion in the reduction algorithm’s termination proof. Neither construction is a counterexample to the principal web-basis conclusion.\par
ID | Severity | Supported conclusion\par
F1 | Major | Proposition 3.21 and the literal monotonicity converse in Corollary 3.33 miss a closed straight trajectory, even in a connected graph.\par
F2 | Major | Algorithm 7.1(4) and Theorem 7.2’s next-iteration reduction claim fail on a graph requiring two preparatory benzene moves.\par
Under the supplied rubric, a demonstrated invalid derivation closely tied to a core result is Major even when a repair is available. This classification does not assert that the main theorem is false. Typographical slips and missing citations alone are outside the finding count. No additional issue has been upgraded merely because an external computation was not rerun.\par
The central claims examined are the graph/tableau bijection of Theorem 6.1, the leading-term result of Theorem 6.4, and the rotation-invariant basis of Theorem 6.10. The quantum statement is over \(\mathbb C(q)\) for \(U_q(\mathrm{sl}_4)\); the semistandard extension in Section 9.3 is at \(q=1\) and \(r=2,3,4\). Rotation invariance is understood in the paper’s stated sense, allowing scalar multiples and the corresponding rotation of the boundary type.\par
Recovery scope. This report completes M25-0120 within the authorized seven-paper recovery. There was no substantive earlier report for this paper to repair; preserved sources and acquisition evidence were reused. Independent scientific acceptance remains pending.\par
\section*{2. F1 — Closed trajectories escape the stated criterion}
Locators. Definitions 3.1 and 3.4 (p. 14); Definitions 3.14, 3.16, and 3.19 (pp. 16–17); Proposition 3.21 (p. 18); Definition 3.30 and Corollary 3.33 (pp. 22–23). These statements were checked in rendered source pages.\par
Take two internal vertices \(v,w\), three distinct parallel simple edges \(a,b,c\) joining them, and boundary edges \(pv,wq\). The clockwise orders are \[v:(p,a,b,c),\qquad w:(q,c,b,a).\] Here \(p\) or \(q\) in a cyclic order denotes the incident boundary edge. Color \(v,q\) black and \(w,p\) white. Every internal vertex has degree four, all multiplicities are one, and the type is \((-1,+1)\). The graph is connected and contracted. Parallel distinct edges are allowed in the ambient category: Remark 3.11 specifically explains why full reducedness subsequently excludes them.\par
(180,124) (90,85)(0,1)30 (90,5)(0,1)30 (90,35)(0,1)50 (90,85)(155,60)(90,35) (90,85)(25,60)(90,35) (90,85)(0,0)\(\bullet\)(90,35)(0,0)\(\circ\) (90,115)(0,0)\(\circ\)(90,5)(0,0)\(\bullet\) (100,111)\(p\)(100,81)\(v\) (100,31)\(w\)(100,1)\(q\) (126,58)\(a\)(94,58)\(b\)(45,58)\(c\)\par
Figure A. The middle edge belongs to the boundary straight strand; the two outer edges form an untested closed straight trajectory.\par
The opposite pairs at \(v\) are \((p,b),(a,c)\), and at \(w\) they are \((q,b),(c,a)\). Thus the only physical boundary \(\operatorname{trip}_2\) strand is \(p-v-b-w-q\), with its reverse representing the same strand. It is simple. The other straight trajectory is the closed walk \(v-a-w-c-v\). Definition 3.16 constructs its trip strands by starting at boundary vertices, so this closed trajectory is not tested.\par
Under \(\varphi\), all internal edges point from \(v\) to \(w\), with \(v\to p\) and \(q\to w\). The internal vertices are respectively a source and a sink, both allowed six-vertex states. Consequently (P1) holds for the sole boundary strand, and (P2) is vacuous. Yet the graph has 2-cycles and fails Definition 3.19 immediately. This disproves the sufficiency in Proposition 3.21. A hypothesis excluding isolated components alone would not remove this example.\par
For the hourglass formulation, the \(\operatorname{trip}_1\) walks are \[\begin{split}
&p-v-c-w-b-v-a-w-q,\\
&q-w-c-v-b-w-a-v-p.
\end{split}\] Each common-vertex set with the boundary \(\operatorname{trip}_2\) strand occupies the entire walk, so it is consecutive along both walks under Definition 3.30’s stated set-based condition. That definition prohibits repeated vertices for \(\operatorname{trip}_2\), but adds no such prohibition for \(\operatorname{trip}_1\). Lemma 3.34 later derives the latter property, so it cannot be silently assumed here.\par
The graph therefore satisfies the written monotonicity predicate. Nevertheless, replace \(a\) by \(v-u-t-w\) with multiplicities \(1,3,1\), using the inverse contraction move of Figure 9. Together with \(b\), this is a forbidden four-cycle. Thus the graph is not fully reduced, contradicting the stated converse of Corollary 3.33.\par
\section*{3. F1 — Core dependency, repair, and bounded checks}
Core dependency and severity. Section 5.8 explicitly replaces the task of proving full reducedness by proving monotonicity (p. 62). The proof of Theorem 5.4 on p. 68 then obtains its full-reducedness assertion from Lemma 5.48 and Corollary 3.33. The missing trajectory condition therefore concerns an advertised structural criterion used in the central proof, rather than an unrelated example. This is the basis for the Major classification.\par
What has and has not been disproved. The counterexample is a legal input to the stated global characterizations. It is not shown to be an output of the growth algorithm. In fact, the growth construction has an additional property that excludes it:\par
> Initially every dangling wire has one original boundary endpoint. An X crosses two wires without joining their straight trajectories, so every surviving dangling wire still has one distinct original endpoint. A cap joins two such wires into a completed boundary-to-boundary trajectory and removes them. Induction therefore excludes every closed straight trajectory in the final linearized diagram.\par
This is a direct check of Algorithm 5.1, not an inference from the faulty global converse. The same argument applies to a suffix diagram used in the growth induction. The correspondence of Theorem 3.25 identifies these straight trajectories with the relevant hourglass trips.\par
A correction should state explicit control of closed straight trajectories or require that every such trajectory have boundary endpoints, prove the corresponding strengthened characterization, and invoke that formulation in Section 5.8. The wire invariant above verifies this extra hypothesis for the growth construction. It removes the particular obstruction exhibited here; this report does not certify an unstated global replacement theorem merely by proposing that hypothesis.\par
\subsection*{Other checks in Sections 2–4}
The labeling and tableau arguments were checked with their actual signed and generalized-rectangle conventions. A generalized rectangle may terminate at \(k(1,1,1,1)\); total signed size need not be zero. Proper-labeling incidence counts give the same signed boundary count for every label, because each internal vertex sees each label once and internal edges cancel. This supplies the balancedness used in Theorem 4.24.\par
Removing a balanced consecutive substring shifts all later prefix vectors by a scalar multiple of \((1,1,1,1)\), leaving adjacent-row slacks unchanged. Concatenating the rectangular lattice words used in the proof likewise preserves the inequalities. Neither argument requires the false assertion that arbitrary nonrectangular lattice words can be freely concatenated.\par
Oscillization uses increasing row order for positive skew columns and removes lower rows first for negative columns. The comparison with promotion uses \(|c_1|\) elementary steps, matching the corresponding number of boundary rotations. The exceptional case in which cutting does not decrease the number of internal vertices is expressly handled by Lemma 4.25 (pp. 35–37).\par
Theorem 2.9’s coefficient formula has a useful noncancellation consequence. For a fixed boundary word, every proper labeling contributes the same sign times a Laurent monomial in \(q\). Thus equality of web invariants cannot hide an earlier boundary word by cancellation; a coefficient \(\pm q^a\) entails a unique contributing labeling. This supports the square-move transfer in Theorem 6.4. These positive checks do not eliminate F1.\par
\section*{4. F2 — A web needing two preparatory benzene moves}
Locators. Definition 3.9 (p. 16); Algorithm 7.1 (pp. 71–72); Figures 43–45 and Theorem 7.2 (pp. 72–73). Step (4) asks for a benzene relation creating a forbidden four-cycle. The proof asserts that a web changed only by that benzene move will undergo Step (2) on the subsequent iteration.\par
Embed three bounded faces in a planar strip with cyclic orders \[\begin{aligned}
Q&=(q_0,q_1,q_2,q_3),\\
H_1&=(q_0,q_1,a_1,b_1,b_0,a_0),\\
H_2&=(b_0,b_1,c_1,d_1,d_0,c_0).
\end{aligned}\] The shared edges are \(q_0q_1\) and \(b_0b_1\); the faces have no further identifications. Give multiplicity two to exactly \[q_0a_0,\quad q_1a_1,\quad b_0c_0,\quad b_1c_1,\quad d_0d_1.\] All remaining internal edges have multiplicity one. Attach a multiplicity-two boundary leaf to each of \(q_2,q_3\), and a simple boundary leaf to each of \(a_0,a_1,c_0,c_1,d_0,d_1\).\par
A bipartition is \[\begin{aligned}
\text{black}&=\{q_0,q_2,a_1,b_0,c_1,d_0\},\\
\text{white}&=\{q_1,q_3,a_0,b_1,c_0,d_1\}.
\end{aligned}\] Each boundary leaf has the opposite color from its neighbor. Every internal vertex has weighted degree four, with exactly one incident doubled edge. The graph is connected and contracted. Boundary order \[(L_{q_2},L_{q_3},L_{a_0},L_{c_0},L_{d_0},L_{d_1},L_{c_1},L_{a_1})\] gives type \((-2,2,1,1,-1,1,-1,-1)\). Splitting each doubled boundary leaf into two consecutive simple leaves gives an oscillating-type example as well.\par
There are twelve internal vertices and fourteen internal edges. The complete list of simple-cycle lengths is \[4,\ 6,\ 6,\ 8,\ 10,\ 12.\] The cycles are \(Q,H_1,H_2\) and the boundaries of the three possible consecutive unions of two or three faces. Thus \(Q\) is the only four-cycle, and \(H_1,H_2\) are the only six-cycles. Boundary leaves create no extra cycle. Initially \(Q\) is entirely simple, so none of Figure 43’s forbidden-square reductions applies; the graph also has no loop to delete under Figure 44.\par
Stage | \(Q\) | \(H_1\) | \(H_2\)\par
Initial | 1111 | 121112 | 121212\par
After flipping \(H_2\) | 1111 | 121212 | 212121\par
After then flipping \(H_1\) | 2111 | 212121 | 112121\par
Here a string lists multiplicities in the displayed cyclic order.\par
Initially only \(H_2\) is benzene-flippable. Its flip doubles \(b_0b_1\), making \(H_1\) flippable, but leaves every edge of \(Q\) simple. Only the second flip, on \(H_1\), doubles \(q_0q_1\) and exposes a forbidden square. The initial web is therefore not fully reduced, while no single available benzene move creates the required four-cycle.\par
\section*{5. F2 — Termination, scope of impact, and other interfaces}
The invalid inference. The principal same-face term of the benzene relation after the first flip in the example still has no forbidden square. It therefore cannot undergo Step (2) on the next iteration. The existence of lower-face side terms in the skein relation does not change that fact about this principal term. Definition 3.9 guarantees a witness somewhere in a move-equivalence class; it does not guarantee a witness after one benzene move.\par
This is a Major defect in the stated reduction procedure and the proof of Theorem 7.2. The alternative basis proof in Theorem 7.3 expressly uses Theorem 7.2 for spanning. The written semistandard spanning argument on p. 80 also invokes Algorithm 7.1. The principal proof in Sections 5–6 does not depend on this immediate-decrease assertion; its separate criterion issue is addressed by F1.\par
What a repair must supply. A revised procedure could select a finite preparatory witness sequence and control both its remaining length and the face count, while treating lower-face skein side terms by an appropriate induction. Such a procedure must justify the permitted preliminary square and contraction moves and prove well-founded progress in every recursive branch. This report has not completed that general replacement proof.\par
Allowing arbitrary benzene moves is insufficient. After the first \(H_2\) flip above, the algorithm could flip \(H_2\) back instead of proceeding to \(H_1\). Proposition 6.7’s directed top-element argument assumes full reducedness and does not give a reduction schedule for this input. Corollary 3.37 also allows sequences of moves, assumes fully reduced graphs, and gives no one-move substitute.\par
An assistant-authored finite graph certificate additionally checks all degrees, bipartition, simple cycles, and multiplicity changes. It executes no author code. The argument above is a complete mathematical certificate of the one-step failure; it does not depend on trusting a saved notebook output.\par
\subsection*{Other checks in Sections 6–9}
The leading terms in Theorem 6.4 are distinct, have coefficients \(\pm q^a\), and hence give linear independence over \(\mathbb C(q)\). The dimension comparison uses rectangular fluctuating tableaux. The basis selection further relies on the directed benzene order and square-move invariance. Reflection need not preserve the top family: the paper explicitly distinguishes top and bottom bases, so this is not a contradiction to its rotation claim.\par
The common-neighbor property used for descents passes between contracted representatives. A vertex sharing two simple boundary leaves can lie on a square move only as a degree-four corner whose remaining edges are square sides; splitting or merging retains the boundary pair on one outer vertex. If the remaining incidence is a single two-hourglass, it is not a square corner. Benzene moves preserve underlying adjacency. This supplies the local justification for the use of Theorem 5.4(v) beyond the initial growth representative in the semistandard argument.\par
In Theorem 9.5, equal boundary arguments at one internal vertex force vanishing by alternation at \(q=1\). Absence of descents within a relabeled block gives a horizontal strip, and standardization recovers the original tableau. Rotation maps weight \(\mu\) to its cyclic rotation; it does not claim to preserve an arbitrary fixed unequal weight.\par
\section*{6. Growth proof and static supplementary audit}
Section 5 was reviewed as an arbitrary-length mathematical argument, with separate attention to what the accompanying notebook actually checks. Applicability, termination, the appliance identities, leading-label uniqueness, and the geometric monotonicity interfaces were followed. Length decreases for caps; the fixed-length strict tilde-lexicographic increase used for the other rules operates on a finite alphabet and supplies the needed well-founded induction.\par
The long-rule argument uses explicit crystal bracketing and finite-state descriptions, rather than an inference from a few computed word lengths. In the state machine of Figure 40, the forbidden-subsequence restrictions are preserved by the relevant raising operators. Once the leading action of \(e_2\) produces the state \(21\langle\overline3,\overline4\rangle2\), no later \(e_2\) action remains. This supports the separation of appliance endpoints used in Lemma 5.48. The additional boundary-coverage invariant needed for F1 is proved on report p. 3.\par
The preserved supplementary notebook is `verifications.ipynb` at commit\par
00002TOKEN\par
in Diagrammatics-Web/sl4-web-basis (https://github.com/Diagrammatics-Web/sl4-web-basis). Its SHA256 is\par
00003TOKEN\par
The historical repository link in reference [GPPSS25a] resolves to this repository. The pinned commit is an acquisition snapshot, not an asserted publication-time release.\par
All 1,357 lines of the static extraction, including every cell source and stored output, were read. No notebook cell or other author code was executed. Saved `True` or `Done.` outputs describe the preserved notebook state and are not reproduced results.\par
Cells | Actual scope visible in the saved notebook\par
14 | All 88 explicitly listed short degeneration rules.\par
16–17 | Two selected long middle words, each with four witnesses.\par
29, 31 | Additional degeneration and identity-action checks only for \(k=0,\ldots,6\).\par
33 | Eight sign lifts of each of ten displayed three-letter patterns: 80 cases.\par
36 | 23 tagged short monotonicity rules: five type-A and eighteen type-B cases; other forms use symmetries.\par
38–39 | Selected long-word monotonicity checks.\par
41–47 | Proper two-strand labelings for six sign patterns and the specified lexicographic conditions.\par
The short-rule degeneration routine compares raising and lowering paths with `zip\_longest` and tests the appliance identity for every pair of those enumerated paths. In the highest-weight implication check, appending two signed letters can reduce an adjacent weight gap by at most two. Truncating each initial nonnegative gap at two therefore explains the 27 canonical gap triples. That cutoff is mathematically motivated. By contrast, the saved long-word examples do not prove the arbitrary-length claims; the manuscript’s inductive arguments remain necessary.\par
No substantive implementation failure was established by this static inspection. The finite author checks have not been independently rerun, and this limitation is retained rather than converted into a finding.\par
\section*{Supporting checks}
The preserved earlier source is arXiv:2306.12501v5 (https://arxiv.org/abs/2306.12501v5), submitted 29 October 2025, with a title-page date of 30 October 2025. It also has 84 pages. Its SHA256 is\par
00004TOKEN\par
The complete token-difference log between the extracted v5 and v6 texts was read: 272 changed blocks, including corrected cross-reference labels, date and pagination changes, and extraction-position differences. The arXiv record describes v6 as fixing a compilation issue. Targeted actual images of v5 pp. 1, 17–18, 22–23, and 72–73 were compared with v6 evidence. The mathematical statements and instructions underlying F1 and F2 persist in both preserved versions. This is not a certification of every graphical detail or of either version’s equivalence to the unavailable journal PDF.\par
\subsection*{Normalization notes excluded from the findings}
Lemma 5.47’s printed appliance direction/index is inconsistent with the direction used in the following monotonicity proof. Reading the crystal definition and row-change argument gives the direction \(\rho_1(a+b)=a+2\) used in that proof. This local transcription issue is not an additional substantive finding.\par
Corollary 6.12’s decreasing row and column orders give distinguished bottom nonzero entries in the coefficient columns; ordinary row-echelon terminology requires adjustment. The triangular/unit-coefficient consequence needed for independence remains available. The isolated promotion-domain and oscillized-index notations on p. 30 were likewise not counted as additional failures: the graph-to-tableau applications use rectangular tableaux and the oscillized boundary size.\par
\subsection*{Actual coverage and verification boundary}
Critical definitions, figures, and proof pages were inspected as rendered images; the exact page lists and image hashes are in the source-inspection record. Small barred labels in growth figures received additional 200-dpi inspection. External foundational theorems were treated as cited dependencies; this review does not purport to reproduce all of their proofs.\par
They are current-conversation collaborators, not independent scientific acceptance. The academic workspace remains Minhao Liu’s Group — Academic Researchers.\par
Only the new recovery staging is written. The original archive, preserved snapshot, and earlier source versions remain read-only. The build uses the already supplied portable compiler in untrusted, cached-only mode, with no author-code execution or shell escape. Build logs bind the current TeX, PDF, compiler, dependencies, and rendered pages.\par
Acceptance state: `scientific\_decision=pending`; `human\_reviews=[]`. Delivery completion does not constitute independent acceptance or author adoption.\par
\end{document}
